%!TEX TS-program = xelatex
%!TEX encoding = UTF-8 Unicode
%
% May 2026 CV — Andrés González Ortega (G1 — InsurTech ES)
% Built on Awesome-CV (posquit0/Awesome-CV, CC BY-SA 4.0)
%

\documentclass[10pt, letterpaper]{awesome-cv}

\geometry{left=1.2cm, top=0.5cm, right=1.2cm, bottom=0.5cm, footskip=.5cm}

% Spacing overrides for 1-page compact fit
\renewcommand{\acvSectionTopSkip}{0.5mm}
\renewcommand{\acvSectionContentTopSkip}{0.3mm}
\renewcommand{\acvHeaderAfterSocialSkip}{1.5mm}
\renewcommand{\acvHeaderAfterNameSkip}{.2mm}
\renewcommand{\acvHeaderAfterPositionSkip}{.2mm}
% Tighter paragraph trailing space
\renewenvironment{cvparagraph}{%
  \vspace{\acvSectionContentTopSkip}%
  \vspace{-3mm}%
  \paragraphstyle%
}{%
  \par%
  \vspace{0.3mm}%
}


% Smaller header name for compact layout
\renewcommand*{\headerfirstnamestyle}[1]{{\fontsize{24pt}{1em}\headerfontlight\color{graytext} #1}}
\renewcommand*{\headerlastnamestyle}[1]{{\fontsize{24pt}{1em}\headerfont\bfseries\color{text} #1}}
\renewcommand*{\headerpositionstyle}[1]{{\fontsize{7.6pt}{1em}\bodyfont\scshape\color{awesome} #1}}
\renewcommand*{\sectionstyleface}[1]{{\fontsize{14pt}{1em}\bodyfont\bfseries #1}}

% Tighter entry spacing
\renewcommand*{\entrytitlestyle}[1]{{\fontsize{9pt}{1em}\bodyfont\bfseries\color{darktext} #1}}
\renewcommand*{\entrypositionstyle}[1]{{\fontsize{8pt}{1em}\bodyfont\scshape\color{graytext} #1}}
\renewcommand*{\entrydatestyle}[1]{{\fontsize{8pt}{1em}\bodyfontlight\slshape\color{graytext} #1}}
\renewcommand*{\entrylocationstyle}[1]{{\fontsize{8pt}{1em}\bodyfontlight\slshape\color{awesome} #1}}
\renewcommand*{\descriptionstyle}[1]{{\fontsize{8.5pt}{1em}\bodyfontlight\upshape\color{text} #1}}
\renewcommand*{\skilltypestyle}[1]{{\fontsize{9pt}{1em}\bodyfont\bfseries\color{darktext} #1}}
\renewcommand*{\skillsetstyle}[1]{{\fontsize{8.5pt}{1em}\bodyfontlight\color{text} #1}}
\renewcommand*{\honortitlestyle}[1]{{\fontsize{8.5pt}{1em}\bodyfont\color{graytext} #1}}
\renewcommand*{\honorpositionstyle}[1]{{\fontsize{8.5pt}{1em}\bodyfont\bfseries\color{darktext} #1}}
\renewcommand*{\honordatestyle}[1]{{\fontsize{8.5pt}{1em}\bodyfont\color{graytext} #1}}
\renewcommand*{\paragraphstyle}{\fontsize{8.5pt}{1.15em}\bodyfontlight\upshape\color{text}}

% PDF metadata
\hypersetup{
  pdftitle={Andrés González Ortega — CV Mayo 2026 G1 InsurTech},
  pdfauthor={Andrés González Ortega},
}

%-------------------------------------------------------------------------------
%  PERSONAL INFORMATION
%-------------------------------------------------------------------------------
\name{Andrés}{González Ortega}
\position{Actuario{\enskip\cdotp\enskip}IA{\enskip\cdotp\enskip}Machine Learning}
\address{Ciudad de México, México}
\mobile{+52 55 4834 4672}
\email{gonorandres@ciencias.unam.mx}
\homepage{gonorandres.github.io}
\github{GonorAndres}

%-------------------------------------------------------------------------------
\begin{document}

\makecvheader[C]

\placetopright{Mayo 2026}{../qr_portfolio.png}

\makecvfooter
  {}
  {\color{lightgray}\textit{\small Andrés Ortega — Pág. \thepage{} de \pageref*{LastPage}}}
  {}

%-------------------------------------------------------------------------------
%  SOBRE MÍ
%-------------------------------------------------------------------------------
\cvsection{Sobre mí}

\begin{cvparagraph}
Soy actuario de la UNAM con \textbf{SOA Exam Probability aprobado}, becario de ciencia de datos en \textbf{Grupo Gigante} y co-autor de artículo científico en el \textbf{IGF-UNAM}. Conecto pensamiento cuantitativo con inteligencia artificial para resolver problemas reales: agentes que navegan miles de artículos regulatorios, modelos que aprenden de millones de registros y plataformas interactivas que acercan la decisión técnica al usuario final. Practico la formación continua y disfruto el trabajo colaborativo con equipos multidisciplinarios. En mi \hrefIcon{https://gonorandres.github.io/blog/}{página web} publico 25+ proyectos en producción, notas técnicas bilingües y aprendizajes abiertos a cualquier disciplina.
\end{cvparagraph}

%-------------------------------------------------------------------------------
%  EXPERIENCIA
%-------------------------------------------------------------------------------
\cvsection{Experiencia}

\begin{cventries}

  \cventry
    {Becario Ciencia de Datos}
    {Grupo Gigante}
    {CDMX}
    {Sept 2025 — Presente}
    {
      \begin{cvitems}
        \item {\textbf{Modelos predictivos} sobre millones de perfiles en seis marcas: \textbf{clasificación} para churn, \textbf{propensión} de compra cruzada y \textbf{segmentación} de cohortes que personalizan ofertas y alimentan campañas de email marketing. Stack: \textbf{Python}, scikit-learn, \textbf{XGBoost}, \textbf{LightGBM}.}
        \item {Contribuyo al \textbf{CDP} del grupo: resolución de identidad entre marcas, pipelines \textbf{ETL} (\textbf{BigQuery}, Cloud SQL, Firestore), dashboards en \textbf{Looker Studio} y modelos de comportamiento web. Iteración directa con marketing, web y operaciones; los dashboards informan decisiones de campaña sobre millones de clientes.}
      \end{cvitems}
    }

  \cventry
    {Asistente de Investigación Científica \normalfont{\textit{(Servicio Social)}}}
    {Instituto de Geofísica, UNAM}
    {CDMX}
    {2025 — Presente}
    {
      \begin{cvitems}
        \item {Co-autor de artículo científico sobre predicción de erupciones del Popocatépetl con \textbf{procesos de renovación Weibull}: modelado estadístico de eventos raros de alto impacto aplicable a riesgo catastrófico y reaseguro. Supervisión: \hrefIcon{https://scholar.google.com/citations?user=lDD0DZQAAAAJ}{Dr.\ Hugo Delgado-Granados.}
      \end{cvitems}
    }

\end{cventries}

%-------------------------------------------------------------------------------
%  HABILIDADES
%-------------------------------------------------------------------------------
\cvsection{Habilidades}

\begin{cvskills}
  \cvskill
    {Lenguajes}
    {Python, R, SQL, TypeScript, Bash, Git, \LaTeX, Excel Avanzado}

  \cvskill
    {IA \& Agentes}
    {LLMs, Claude API, Claude Agent SDK, RAG, FTS5/BM25, Transformers, Inferencia Bayesiana, NLP}

  \cvskill
    {Ciencia Actuarial}
    {Seguros Vida/Daños/GMM, Tarificación (GLM, frecuencia-severidad, credibilidad Bühlmann-Straub), Reservas (BEL, IBNR, Chain Ladder, BF), Lee-Carter, Whittaker-Henderson, RCS/SCR bajo LISF/CUSF (análogo a Solvency II), EMSSA-09, Reaseguro, CONAPO/HMD, Pensiones IMSS (Ley 73/97, Fondo Bienestar)}

  \cvskill
    {Cloud \& DevOps}
    {GCP (Cloud Run, BigQuery, Pub/Sub), Docker, Terraform, Dagster, GitHub Actions, PostgreSQL}

  \cvskill
    {DS \& ML}
    {scikit-learn, PyTorch, XGBoost, LightGBM, Monte Carlo, SHAP, FastAPI, Streamlit, Power BI, Looker Studio}

  \cvskill
    {Idiomas}
    {Inglés avanzado (lectura, redacción, comunicación fluida)}
\end{cvskills}

%-------------------------------------------------------------------------------
%  EDUCACIÓN
%-------------------------------------------------------------------------------
\cvsection{Educación}

\begin{cventries}
  \cventry
    {Licenciatura en Actuaría — Pasante (100\% de créditos cubiertos)}
    {UNAM, Facultad de Ciencias}
    {}
    {2021 — 2025}
    {
      \begin{cvitems}
        \item {Titulación por exámenes profesionales SOA (\hrefIcon{https://www.soa.org/}{Society of Actuaries}).}
      \end{cvitems}
    }
\end{cventries}

%-------------------------------------------------------------------------------
%  CERTIFICACIONES
%-------------------------------------------------------------------------------
\cvsection{Certificaciones}

\begin{cvhonors}
  \cvhonor
    {SOA Exam Probability}
    {Society of Actuaries — \hrefIcon{https://drive.google.com/file/d/1rt3emgBnPpQi7NiXkBQeA5WwJTV0JuAf/view?usp=sharing}{Aprobado}}
    {}
    {Mar 2026}

  \cvhonor
    {Databricks Fundamentals}
    {\hrefIcon{https://credentials.databricks.com/45f31d0f-ca80-4e3a-80c1-ede89826f6ce}{Databricks Academy}}
    {}
    {Abr 2026}

  \cvhonor
    {Data Scientist Professional with R}
    {\hrefIcon{https://drive.google.com/file/d/1Woe6xqxofloFZ9uM2f5VsdGxY-WQWCRI/view?usp=drive_link}{DataCamp}}
    {}
    {2024}

  \cvhonor
    {Associate Data Scientist in R}
    {\hrefIcon{https://drive.google.com/file/d/1xVLk15A1LBbyXH1fZnd4eBTaJ4d8FLLO/view?usp=drive_link}{DataCamp}}
    {}
    {2024}

  \cvhonor
    {Associate Data Analyst}
    {\hrefIcon{https://drive.google.com/drive/folders/1-QUuVH3zPBPsl6RhUvuUWxpZkP1i8TYb?usp=sharing}{DataCamp}}
    {}
    {2024}
\end{cvhonors}

%-------------------------------------------------------------------------------
%  PROYECTOS
%-------------------------------------------------------------------------------
\cvsection{Proyectos}

\begin{cventries}

  \cventry
    {Asistente de Regulación Actuarial (RAG)}
    {}
    {}
    {\hrefIcon{https://actuarial-regulation-agent-d3qj5vwxtq-uc.a.run.app/}{app live}}
    {
      \begin{cvitems}
        \item {\hl{2,354} artículos de LISF y CUSF indexados con \hl{2,882} referencias cruzadas como grafo de conocimiento. Búsqueda \hl{FTS5/BM25} con pesos por columna en milisegundos; responde con cita exacta. Despliegue dual: Claude API + Qwen 72B open-source. \textit{Claude Agent SDK, FastAPI, GCP.}}
      \end{cvitems}
    }

  \cventry
    {SIMA — Sistema Integral de Modelación Actuarial}
    {}
    {}
    {\hrefIcon{https://sima-d3qj5vwxtq-uc.a.run.app}{app live}}
    {
      \begin{cvitems}
        \item {Valuar reservas de vida exige conectar mortalidad, producto y capital en un solo flujo. Integra \hl{Lee-Carter} (\hl{77.7\%} varianza explicada, México 1990-2019), conmutación, \hl{BEL} para tres productos y \hl{RCS} bajo \hl{LISF} con pruebas de estrés. Impacto pandemia cuantificado (+3-10\% prima). 240 tests automatizados. \textit{Python, FastAPI, React, GCP.}}
      \end{cvitems}
    }

  \cventry
    {GMM Explorer — Tarificador de Gastos Médicos}
    {}
    {}
    {\hrefIcon{https://gmm-explorer.vercel.app/contexto}{app live}}
    {
      \begin{cvitems}
        \item {Tarificar GMM sin datos reales es especular. Procesa \hl{5.1M} siniestros de la CNSF (2020-2024), los clasifica en tres niveles con IA y calcula prima pura por \hl{frecuencia-severidad} ajustada por inflación médica. \textit{Python, Next.js, Vercel.}}
      \end{cvitems}
    }

  \cventry
    {CreditGraph — Riesgo Crediticio Topológico}
    {}
    {}
    {\hrefIcon{https://gonorandres.github.io/blog/credit-graph-topological-risk/}{blog}}
    {
      \begin{cvitems}
        \item {SQL trata cada préstamo como evento independiente; garantías circulares crean exposición invisible. Grafo en \hl{Neo4j} con \hl{PySpark} en Databricks, scoring \hl{LightGBM}+Platt calibrado. \textit{Python, Cypher, Databricks.}}
      \end{cvitems}
    }

\end{cventries}

\end{document}
